% =====================================================================
%  VR 검도훈련 콘텐츠 구성요소 — 델파이 설문 문항 (단일 원본)
%  설문지(survey/delphi-survey.tex)와 논문 부록이 모두 이 파일을 사용합니다.
%
%   \Domain{번호}{영역명}{영역 설명}      대분류
%   \Module{번호}{중분류명}                중분류 (표 안의 구분 행)
%   \Item{번호}{구성요소명}{설명}          세부 구성요소 (평정 문항)
%   \NewItem{번호}{구성요소명}{설명}       연구자 추가 제안 문항 ([신규] 표시,
%                                          \ShowProposedItems 가 false 면 숨김)
%   \EndDomain                            영역 종료 (영역 타당성 + 개방형 의견란)
% =====================================================================

% ---------------------------------------------------------------------
\Domain{1}{하드웨어 및 인터랙션 (Hardware \& Interaction)}
  {가상 공간에서 실제 검도의 신체 감각, 칼의 물리적 궤적 및 반동을 구현하기 위한 기술적·장비적 요소}

\Module{1.1}{모션 트래킹}
\Item{1.1.1}{9축/6DoF 죽도 트래킹}{죽도의 위치, 속도, 칼날 각도를 실시간 추적}
\Item{1.1.2}{HMD 시선 추적}{사용자의 시선 방향 및 주시점 인식}
\Item{1.1.3}{하체 발놀림(Foot) 추적}{스텝, 발밟기 등 하체 동작 트래킹}

\Module{1.2}{햅틱 피드백}
\Item{1.2.1}{타격 순간 반동 피드백}{격자 부위 타격 시 죽도 장비에 물리적 진동/충격 전달}
\Item{1.2.2}{교검(交劍) 저항 피드백}{상대 죽도와 칼이 맞닿거나 쳐누를 때 물리적 저항감 제공}

\Module{1.3}{공간 음향}
\Item{1.3.1}{3D 입체 음향}{타격음, 상대 기합 소리, 경기장/도장 울림 등 공간 음향 구현}
\EndDomain

% ---------------------------------------------------------------------
\Domain{2}{검도 무도론 및 교수학습 (Kendo Pedagogy)}
  {검도 고유의 수련 체계, 거리감(간합), 무도적 가치를 VR 환경에 맞춰 구조화한 콘텐츠 요소}

\Module{2.1}{수련 단계 모드}
\Item{2.1.1}{기본 수련}{중단자세, 기본 발놀림, 공간치기 반복 수련 모드}
\Item{2.1.2}{응용/격권 수련}{머리·손목·허리 연속 치기 및 받아치기/쳐내기 기법 모드}
\Item{2.1.3}{자유 대련}{AI 또는 멀티플레이어 기반 대적 시뮬레이션 모드}

\Module{2.2}{간합 및 AI}
\Item{2.2.1}{일간의 거리(一間之間) 시각화}{유효 타격이 가능한 시공간적 거리 안내 표시}
\Item{2.2.2}{대적 AI 엔진}{유저의 거리 조절에 맞춰 전진/후진 및 응대 반격을 수행하는 AI}

\Module{2.3}{기세 및 기합}
\Item{2.3.1}{기합 음성 인지}{마이크 입력을 통해 타격 순간의 기합 소리 크기 및 타이밍 인식}

\NewModule{2.4}{예법 및 무도 가치}
\NewItem{2.4.1}{예법(禮法) 수련}{입장·례(禮)·준거(蹲踞)·납도 등 검도 예법 절차 안내 및 수행 여부 확인}
\EndDomain

% ---------------------------------------------------------------------
\Domain{3}{동작 평가 및 피드백 (Assessment \& Feedback)}
  {수련자의 동작이 검도의 정격(正擊) 규칙 및 원리에 부합하는지 측정·평가하는 요소}

\Module{3.1}{심검체 일치도}
\Item{3.1.1}{심검체(心劍體) 동기화 측정}{시선(心)-죽도 궤적(劍)-발밟기(體)의 타임스탬프 일치 판정}

\Module{3.2}{정격 판정}
\Item{3.2.1}{칼날 각도(刃筋) 검증}{타격 시 죽도 날지름이 유효 범위(45도 이내)로 들어왔는지 판정}
\Item{3.2.2}{유효 격자부위 판정}{머리, 손목, 허리, 찌르기 부위에 정확히 타격되었는지 측정}
\NewItem{3.2.3}{타돌부(物打) 판정}{죽도의 유효 타돌부(선단 약 1/4 부위)로 타격했는지 판정}
\NewItem{3.2.4}{잔심(殘心) 판정}{타격 후 자세·기세·거리를 유지하며 다음 동작에 대비했는지 판정}

\Module{3.3}{다차원 피드백}
\Item{3.3.1}{3D 슬로우 모션 리플레이}{타격 순간을 360도 및 속도 조절로 재확인하는 기능}
\Item{3.3.2}{수련 데이터 시각화}{타격 속도(m/s), 반응 시간, 죽도 궤적 흔들림 다이어그램 제공}
\EndDomain

% ---------------------------------------------------------------------
\Domain{4}{시스템 및 환경 (UI/UX \& System)}
  {훈련자의 몰입감 유지, 개인 맞춤화 및 연속적인 학습 관리를 지원하는 보조 요소}

\Module{4.1}{수련 환경 구축}
\Item{4.1.1}{가상 훈련 공간 선택}{기본 도장, 체육관, 전국 대회 경기장 등 가상 환경 제공}

\Module{4.2}{아바타 커스텀}
\Item{4.2.1}{신체 치수 보정}{사용자의 키, 팔 길이, 죽도 규격에 맞춘 1:1 아바타 보정}

\Module{4.3}{학습 관리(LMS)}
\Item{4.3.1}{수련 리포트 제공}{누적 수련 시간, 부위별 타격 성공률, 심검체 일치도 추이 그래프 제공}
\NewItem{4.3.2}{지도자 모니터링 및 원격 코칭}{지도자가 수련자의 기록을 열람하고 코멘트·과제를 제공하는 기능}

\NewModule{4.4}{안전 및 사용 편의}
\NewItem{4.4.1}{안전 경계 및 충돌 방지}{실제 공간의 안전 경계(가디언) 설정 및 죽도 휘두르기 공간 확보 안내}
\NewItem{4.4.2}{사이버멀미 및 피로 관리}{프레임률 유지, 이동 방식 선택, 휴식 알림 등 VR 멀미·신체 피로 관리}
\EndDomain

% ---------------------------------------------------------------------
\Domain{5}{게임 설계 및 몰입 (Game Design \& Engagement)}
  {수련자의 지속적인 훈련 동기 유발, 몰입(Flow) 유지 및 학습 효과 극대화를 위한 게임화 요소}

\Module{5.1}{동기부여 및 보상}
\Item{5.1.1}{단/급 승단 체계}{수련 성과 기반 단/급 수여 및 진급 칭호 제공}
\Item{5.1.2}{수련 미션/퀘스트}{일일 수련 목표 설정 및 미션 달성 시 보상 부여}

\Module{5.2}{난이도 및 밸런싱}
\Item{5.2.1}{동적 난이도 조절(DDA)}{수련자 숙련도 및 정격률에 맞춘 AI 공격 속도/패턴 실시간 조절}
\Item{5.2.2}{수련 커브 디자인}{입문부터 고단자 수련까지 난이도 상승 곡선(Learning Curve) 단계화}

\Module{5.3}{몰입 및 연출}
\Item{5.3.1}{정격 타격 이펙트}{무도적 품위를 유지하는 수준의 정격 시각/청각 피드백 연출}

\Module{5.4}{경쟁 및 비동기 수련}
\Item{5.4.1}{랭킹 및 리더보드}{정격 정확도 및 심검체 일치도 점수 기반 랭킹 시스템}
\Item{5.4.2}{고단자 고스트(Ghost) 대전}{고단자의 실제 모션 데이터를 기반으로 한 비동기 대적 수련}
\EndDomain
